% Section 1.6, translated from sv/chapters/01/summary.tex.

\section{Summary}
\label{c1:sec:sammanfattning}

\begin{booktable}{@{}p{15em}L@{}}
  \thead{Concept} & \thead{Definition or formula} \\
  \midrule
  Absolute value & $\abs{a} \geq 0$, the distance to zero \\
  Order of operations & Brackets, powers, $\times$ and $\div$, $+$ and $-$ \\
  Adding fractions & Common denominator \\
  Multiplying fractions & Numerator $\times$ numerator, denominator $\times$ denominator \\
  Percentage & $p\,\% = \frac{p}{100}$ \\
  Parallel resistance & $\frac{1}{R_{\text{TOT}}} = \frac{1}{R_1} + \frac{1}{R_2}$ \\
  Parallel resistance, two resistors & $R_{\text{p}} = \frac{R_1 \times R_2}{R_1 + R_2}$ \\
  Series resistance & $R_{\text{s}} = R_1 + R_2$ \\
  Ohm's law & $U = R \times I$ \\
  Efficiency & $\eta = \frac{P_{\text{out}}}{P_{\text{in}}} \times 100\,\%$ \\
\end{booktable}

\needspace{6\baselineskip}
\subsection*{What you should be able to do}
After this chapter, you should be able to:
\begin{itemize}
  \item Know the different number sets and their properties.
  \item Apply the order of operations and the basic arithmetic operations.
  \item Calculate with fractions, decimals and percentages.
  \item Calculate resistance, current and voltage in simple series and parallel connections.
\end{itemize}

\testadigsjalv
\begin{enumerate}
  \item Which number set does $\sqrt{5}$ belong to? Justify your answer briefly.
  \item Calculate the parallel resistance of $R_1 = 4\,\Omega$ and $R_2 = 12\,\Omega$.
\end{enumerate}

\needspace{6\baselineskip}
\subsection*{Next chapter}
\Kapref{2} is about algebra: algebraic expressions, simplification and factorisation.
